\section{One Hilbert space behind all cylinder depths}
\label{a-limit-hilbert}

An isometric refinement map compares two finite registers. An infinite
quantum dynamics requires more: a common completed Hilbert space, a
specified operator domain, and a topology in which the cutoffs converge.
These requirements can all be met for the cylinder tower. They are separate
requirements for a walk on a growing finite tree.

Fix a prime $p$. The compact additive group $\Z_p$ has normalized Haar
measure $\mu_p$, with
\[
 B_{m,a}=a+p^m\Z_p,\qquad 0\leq a<p^m,\qquad
 \mu_p(B_{m,a})=p^{-m}.
\]
Write $\HH_p=L^2(\Z_p,\mu_p;\C)$. We use the complex inner product
$\inner{f}{g}=\int\overline f g\,d\mu_p$, linear in its second argument.
The $p$-adic coordinate belongs to the measurable space. Amplitudes and the
Born probabilities use ordinary complex Hilbert space.

\subsection{Counting amplitudes and Haar densities}

At depth $m$, the finite space $\HH_{p,m}=\ell^2(\Z/p^m\Z)$ has counting
inner product and orthonormal basis $e_a$. Define
\begin{equation}
 \mathcal U_{p,m}e_a=p^{m/2}\mathbf1_{B_{m,a}},
 \qquad \mathcal V_m=\operatorname{span}
 \{\mathbf1_{B_{m,a}}:0\leq a<p^m\}.
 \label{a-limit-cylinder-embedding}
\end{equation}
The cylinders are disjoint. Consequently
\[
 \inner{\mathcal U_{p,m}e_a}{\mathcal U_{p,m}e_b}
 =p^m\mu_p(B_{m,a}\cap B_{m,b})=\delta_{ab}.
\]
The same indicators span $\mathcal V_m$, so $\mathcal U_{p,m}$ is unitary
onto this subspace. For $c=\sum_a c_ae_a$, its image has constant value
$p^{m/2}c_a$ on $B_{m,a}$. In particular,
\begin{equation}
 \int_{B_{m,a}}|\mathcal U_{p,m}c|^2\,d\mu_p=|c_a|^2.
 \label{a-limit-born-normalization}
\end{equation}
There is no extra probability factor after this normalization.

The finite refinement map is
\[
 V_{p,m}e_a=\frac1{\sqrt p}
 \sum_{d=0}^{p-1}e_{a+d p^m}.
\]
Since $B_{m,a}=\bigsqcup_d B_{m+1,a+d p^m}$, it satisfies
\begin{equation}
 \mathcal U_{p,m+1}V_{p,m}=\mathcal U_{p,m}.
 \label{a-limit-common-embedding}
\end{equation}
Indeed, the coefficient on each child indicator is
$p^{-1/2}p^{(m+1)/2}=p^{m/2}$, precisely the coefficient on the parent.
Thus the refinement changes the address resolution of an embedded state
without changing the function in $\HH_p$.

An equivalent finite convention uses the weighted norm
\[
 \norm{f}_{m,\mathrm{Haar}}^2=p^{-m}\sum_a|f(a)|^2.
\]
Here $f(a)$ is a density amplitude rather than a counting amplitude, and
the embedding into $\HH_p$ sends $f$ to
$\sum_a f(a)\mathbf1_{B_{m,a}}$. Its refinement repeats the parent value
on all children, without $p^{-1/2}$. The coordinate change
$f(a)=p^{m/2}c_a$ identifies the two conventions. The factors
$p^{-m}$ in the measure and $p^{m/2}$ in the counting embedding are a
single normalization expressed in two places.

\subsection{Cylinder averaging as an orthogonal projection}

For $f\in\HH_p$, define
\begin{equation}
 (\mathsf E_m f)(x)
 =p^m\int_{x+p^m\Z_p} f(y)\,d\mu_p(y).
 \label{a-limit-expectation}
\end{equation}
The integral is defined since $L^2$ is contained in $L^1$ on this
probability space. It produces a function constant on each depth-$m$
cylinder.

\begin{proposition}[Nested cylinder projections]
\label{a-limit-expectation-algebra}
The operator $\mathsf E_m$ is the orthogonal projection onto $\mathcal V_m$.
For $m,n\geq0$,
\begin{equation}
 \mathsf E_m^*=\mathsf E_m,\qquad
 \mathsf E_m^2=\mathsf E_m,\qquad
 \mathsf E_m\mathsf E_n=\mathsf E_{\min(m,n)}.
 \label{a-limit-projection-product}
\end{equation}
Moreover $\mathsf E_m f\to f$ in $L^2$ for every $f\in\HH_p$.
\end{proposition}
\begin{proof}
For a cylinder-constant $g$, the integral of $f-\mathsf E_m f$ over each
depth-$m$ cylinder is zero. Multiplying by the constant value of
$\overline g$ there and summing shows
$\inner{g}{f-\mathsf E_mf}=0$. The image is in $\mathcal V_m$, and averaging
a function in $\mathcal V_m$ returns the function. These statements
characterize the orthogonal projection and establish the first two
identities. The spaces are nested, because every cylinder partitions into
its children. Orthogonal projections onto nested subspaces satisfy the
third identity: the larger projection fixes the smaller range, and the
reverse order follows by taking adjoints.

For completeness, the union of these ranges is dense. Cylinders form a
countable clopen base for the topology on $\Z_p$ and generate its Borel
sigma-algebra. If $f\in L^2$ is orthogonal to every cylinder indicator,
the finite complex measure $\nu(A)=\int_A f\,d\mu_p$ vanishes on
cylinders and their finite unions. The class of sets on which $\nu$ is
zero contains the whole space, is closed under complements, and is
closed under countable disjoint unions. The elementary monotone-class
extension from the cylinder algebra therefore gives $\nu(A)=0$ for every
Borel set. Applying this to the positive and negative level sets of the
real and imaginary parts of $f$ yields $f=0$ almost everywhere.
The orthogonal complement of $\bigcup_m\mathcal V_m$ is thus zero.

Given $\varepsilon>0$, choose $g\in\mathcal V_M$ with
$\norm{f-g}_2<\varepsilon$. For $m\geq M$, the projection fixes $g$, so
$\norm{f-\mathsf E_m f}_2
=\norm{(I-\mathsf E_m)(f-g)}_2\leq\norm{f-g}_2<\varepsilon$.
This proves strong convergence without exchanging an infinite integral
with a limit.
\end{proof}

The completion of the isometric direct system
$(\HH_{p,m},V_{p,m})$ is therefore $\HH_p$: equation
\eqref{a-limit-common-embedding} gives a single isometric map on the
algebraic direct limit, and its image is the dense union of the
$\mathcal V_m$. For $m\leq M$, with
$V_{p,m\to M}=V_{p,M-1}\cdots V_{p,m}$, one also has
\begin{equation}
 \mathsf E_m\big|_{\mathcal V_M}
 =\mathcal U_{p,M}V_{p,m\to M}V_{p,m\to M}^*
 \mathcal U_{p,M}^*.
 \label{a-limit-finite-projection}
\end{equation}
Both sides are the orthogonal projection in $\mathcal V_M$ onto
$\mathcal V_m$. This is the exact relation between the finite matrices
and the common infinite Hilbert space.

\subsection{Orthogonal details and a complete Haar basis}

Use $\mathsf F_0=\mathsf E_0$ and
$\mathsf F_m=\mathsf E_m-\mathsf E_{m-1}$ for $m\geq1$. These are the
infinite depth details called $\mathsf D_m$ in the archived chapter.
They correspond to the finite $D_m^{(M)}$ after the embedding
$\mathcal U_{p,M}$. Equation \eqref{a-limit-projection-product} gives
\begin{equation}
 \mathsf F_m\mathsf F_n=\delta_{mn}\mathsf F_m,
 \qquad \sum_{m=0}^M\mathsf F_m=\mathsf E_M.
 \label{a-limit-detail-algebra}
\end{equation}
For example, if $1\leq m<n$, the four products in
$(\mathsf E_m-\mathsf E_{m-1})(\mathsf E_n-\mathsf E_{n-1})$
are $\mathsf E_m,\mathsf E_m,\mathsf E_{m-1},\mathsf E_{m-1}$ and cancel.
For $m=n$ the same expansion is
$\mathsf E_m-\mathsf E_{m-1}$. The cases involving $m=0$ use
$\mathsf E_0\mathsf E_n=\mathsf E_0$.

Consequently
\begin{equation}
 \HH_p=\bigoplus_{m=0}^{\infty}\operatorname{Ran}\mathsf F_m,
 \qquad \norm{f}_2^2=\sum_{m=0}^{\infty}\norm{\mathsf F_mf}_2^2.
 \label{a-limit-parseval}
\end{equation}
The sum of the projections is $I$ in the strong operator topology.
The detail ranks are $d_0=1$ and $d_m=(p-1)p^{m-1}$ for $m\geq1$,
the differences of the dimensions $p^m-p^{m-1}$.

For an explicit basis, put $\omega_p=\exp(2\pi i/p)$ and define
\begin{equation}
 \psi_{m,a,q}=p^{(m-1)/2}\sum_{d=0}^{p-1}
 \omega_p^{qd}\mathbf1_{B_{m,a+d p^{m-1}}},
 \quad
 \begin{matrix}m\geq1,\quad 0\leq a<p^{m-1},\\1\leq q<p.\end{matrix}
 \label{a-limit-haar-wavelets}
\end{equation}
Each has zero average over its parent, since
$\sum_d\omega_p^{qd}=0$, and lies in $\operatorname{Ran}\mathsf F_m$.
Its norm squared is $p^{m-1}p\,p^{-m}=1$. Distinct parents have disjoint
support, while two characters at the same parent have inner product
$p^{-1}\sum_d\omega_p^{(q'-q)d}=\delta_{qq'}$. The number of these
orthonormal functions is exactly $d_m$, so they span that level.
Together with $\psi_0=1$, they form a complete orthonormal basis by
\eqref{a-limit-parseval}.

\subsection{A statewise approximation rate}

For $s>0$, define the Haar regularity norm
\[
 \norm{f}_{\mathfrak H_p^s}^2
 =\sum_{m\geq0}p^{2sm}\norm{\mathsf F_mf}_2^2,
 \qquad
 \mathfrak H_p^s=\{f\in\HH_p:\norm{f}_{\mathfrak H_p^s}<\infty\}.
\]
No identification with another Sobolev convention is required here.
The definition itself gives the quantitative cutoff estimate.

\begin{proposition}[Exact tail and regularity bound]
\label{a-limit-tail}
For $f\in\HH_p$ and every $M\geq0$,
\begin{equation}
 \norm{(I-\mathsf E_M)f}_2^2
 =\sum_{m>M}\norm{\mathsf F_mf}_2^2.
 \label{a-limit-tail-identity}
\end{equation}
If $f\in\mathfrak H_p^s$, then
\begin{equation}
 \norm{(I-\mathsf E_M)f}_2
 \leq p^{-s(M+1)}\norm{f}_{\mathfrak H_p^s}.
 \label{a-limit-tail-bound}
\end{equation}
The convergence $\mathsf E_M\to I$ is strong and is not convergence in
operator norm.
\end{proposition}
\begin{proof}
The omitted details are orthogonal, so Parseval gives the equality.
For $m>M$, $p^{2sm}\geq p^{2s(M+1)}$. Multiplying the omitted summands
by this lower bound proves the squared estimate, and taking square roots
gives \eqref{a-limit-tail-bound}. At every finite $M$ there is a unit
vector in $\operatorname{Ran}\mathsf F_{M+1}$, on which
$I-\mathsf E_M$ acts as the identity. Thus
$\norm{I-\mathsf E_M}=1$ for every $M$.
\end{proof}

The rate is a bound on states in the displayed regularity class. For an
arbitrary $L^2$ state the tail still tends to zero, but no common geometric
rate follows. A finite array already in $\mathcal V_L$ has zero tail once
$M\geq L$, which is a different statement from testing an infinite tail.


\section{Self-adjoint dynamics and what a cutoff converges to}
\label{a-limit-dynamics}

\subsection{The maximal domain for level energies}

Let $(\epsilon_m)_{m\geq0}$ be any sequence of real numbers. Define
\begin{align}
 \Dom(H_\epsilon)
 &=\left\{f\in\HH_p:
 \sum_{m\geq0}|\epsilon_m|^2\norm{\mathsf F_mf}_2^2<\infty\right\},
 \nonumber\\
 H_\epsilon f&=\sum_{m\geq0}\epsilon_m\mathsf F_mf.
 \label{a-limit-hamiltonian-domain}
\end{align}
The series converges in $L^2$ exactly on this domain. It is not necessary
to bound the energies uniformly, but it is necessary to state the domain.

\begin{theorem}[Real level energies define a self-adjoint operator]
\label{a-limit-selfadjoint}
The operator in \eqref{a-limit-hamiltonian-domain} is densely defined,
closed, and self-adjoint. The finite cylinder functions form an operator
core. It is bounded if and only if $\sup_m|\epsilon_m|<\infty$, and its
spectrum is
\begin{equation}
 \sigma(H_\epsilon)=\overline{\{\epsilon_m:m\geq0\}}.
 \label{a-limit-spectrum}
\end{equation}
\end{theorem}
\begin{proof}
The orthogonal decomposition \eqref{a-limit-parseval} identifies $f$ with
the square-summable sequence $f_m=\mathsf F_mf$. Finite sequences belong
to the domain, so it is dense. The real multipliers imply symmetry.
To identify the adjoint, suppose $g\in\Dom(H_\epsilon^*)$ and write
$k=H_\epsilon^*g$. Testing the adjoint identity against an arbitrary
vector supported on level $m$ gives $k_m=\epsilon_m g_m$. Since $k$ is
square summable, $\sum_m|\epsilon_m|^2\norm{g_m}^2<\infty$.
Thus $g\in\Dom(H_\epsilon)$ and the adjoint agrees with $H_\epsilon$.
The reverse inclusion is symmetry. This proves self-adjointness and
therefore closedness.

For $f\in\Dom(H_\epsilon)$, the cylinder truncation $\mathsf E_M f$
satisfies both
$\norm{f-\mathsf E_M f}_2\to0$ and
$\norm{H_\epsilon(f-\mathsf E_M f)}_2\to0$, because the two squared norms
are the tails of convergent nonnegative series. This is precisely density
in the graph norm and proves the core assertion.
If the energies have supremum $C<\infty$, Parseval gives
$\norm{H_\epsilon f}_2\leq C\norm f_2$. Conversely, a unit vector on
level $m$ has image norm $|\epsilon_m|$, so boundedness forces a finite
supremum.

If $z$ is outside the closed energy set, the sequence
$(\epsilon_m-z)^{-1}$ is bounded and defines an inverse for
$H_\epsilon-z$. Its image lies in the domain because
$\epsilon_m/(\epsilon_m-z)=1+z/(\epsilon_m-z)$ is bounded as well.
This proves that such $z$ is in the resolvent set. An energy
$\epsilon_m$ is an eigenvalue, while a limit point $z$ has unit level
vectors $u_j$ with $\norm{(H_\epsilon-z)u_j}_2\to0$.
A bounded inverse would contradict this limit. Hence the closed energy
set is exactly the spectrum.
\end{proof}

This proof also excludes a hidden choice of boundary condition at infinite
depth. On the finite-cylinder core, the real level operator has the unique
self-adjoint closure just described.

\subsection{Explicit spectral calculus and strong continuity}

For bounded Borel $F:\R\to\C$, define
\begin{equation}
 F(H_\epsilon)f=\sum_{m\geq0}F(\epsilon_m)\mathsf F_mf.
 \label{a-limit-borel-calculus}
\end{equation}
The series converges by Parseval and
$\norm{F(H_\epsilon)f}_2\leq\sup_m|F(\epsilon_m)|\norm f_2$.
The same level decomposition verifies sums, products, and adjoints of
these multipliers directly. In particular,
\begin{align}
 U(t)&=e^{-itH_\epsilon}
 =\sum_{m\geq0}e^{-it\epsilon_m}\mathsf F_m,
 \label{a-limit-unitary}\\
 (z-H_\epsilon)^{-1}
 &=\sum_{m\geq0}(z-\epsilon_m)^{-1}\mathsf F_m,
 \quad z\notin\sigma(H_\epsilon).
 \label{a-limit-resolvent}
\end{align}
The projections in these sums converge strongly, not absolutely in
operator norm in general.

The modulus-one multipliers give $U(t)^*U(t)=I$ and
$U(t+s)=U(t)U(s)$. They also give strong continuity:
\[
 \norm{(U(t)-U(t_0))f}_2^2
 =\sum_m|e^{-it\epsilon_m}-e^{-it_0\epsilon_m}|^2
 \norm{\mathsf F_mf}_2^2\longrightarrow0.
\]
To justify this limit, first bound the omitted summands by
$4\sum_{m>M}\norm{\mathsf F_mf}_2^2$, then use ordinary continuity in
the finite sum up to $M$. On $\Dom(H_\epsilon)$ the same argument,
using $|(e^{-it\epsilon}-1)/t|\leq|\epsilon|$, proves
\[
 \lim_{t\to0}\frac{U(t)f-f}{t}=-iH_\epsilon f.
\]
Thus the declared self-adjoint operator is the generator. A general
$L^2$ state evolves unitarily even if this derivative does not exist.

\subsection{Two distinct extensions of the finite dynamics}

Let $H_\epsilon^{(M)}$ be the finite restriction of $H_\epsilon$ to
$\mathcal V_M$. One useful approximation to $U(t)$ is its embedded
finite propagator, extended by zero on the omitted subspace:
\begin{equation}
 U_M^{\mathrm{emb}}(t)
 =\sum_{m=0}^{M}e^{-it\epsilon_m}\mathsf F_m
 =U(t)\mathsf E_M.
 \label{a-limit-embedded-propagator}
\end{equation}
It is unitary on $\mathcal V_M$ and has kernel
$\mathcal V_M^\perp$ as an operator on $\HH_p$.

\Needspace{15\baselineskip}
\begin{proposition}[Exact state error for a compatible cutoff]
\label{a-limit-unitary-error}
For every $f\in\HH_p$,
\begin{equation}
 \norm{U(t)f-U_M^{\mathrm{emb}}(t)f}_2
 =\norm{(I-\mathsf E_M)f}_2.
 \label{a-limit-unitary-tail}
\end{equation}
This equality is uniform for all real $t$. For
$f\in\mathfrak H_p^s$, the right side is bounded by
$p^{-s(M+1)}\norm f_{\mathfrak H_p^s}$.
For $f\in\Dom(H_\epsilon)$ there is also the exact graph-norm identity
\begin{align}
 &\norm{U(t)f-U_M^{\mathrm{emb}}(t)f}_2^2
 +\norm{H_\epsilon(U(t)f-U_M^{\mathrm{emb}}(t)f)}_2^2
 \nonumber\\
 &\hspace{2em}=
 \sum_{m>M}(1+|\epsilon_m|^2)\norm{\mathsf F_mf}_2^2.
 \label{a-limit-graph-tail}
\end{align}
\end{proposition}
\begin{proof}
The difference is the omitted orthogonal sum
$\sum_{m>M}e^{-it\epsilon_m}\mathsf F_mf$. Its multipliers have modulus
one, so Parseval gives \eqref{a-limit-unitary-tail}. Applying
$H_\epsilon$ multiplies the omitted level by $\epsilon_m$ and gives the
second identity. The domain hypothesis ensures that this second series
is finite. The regularity estimate is Proposition \ref{a-limit-tail}.
\end{proof}

An alternative is to extend the finite Hamiltonian by the zero
Hamiltonian on the omitted subspace:
\[
 \widetilde H_M=H_\epsilon\mathsf E_M
 =\sum_{m\leq M}\epsilon_m\mathsf F_m.
\]
This is a bounded self-adjoint operator on all of $\HH_p$. Its
propagator is a full-space unitary, but has a different tail:
\begin{equation}
 e^{-it\widetilde H_M}
 =U_M^{\mathrm{emb}}(t)+(I-\mathsf E_M),
 \qquad
 \norm{(e^{-it\widetilde H_M}-U(t))f}_2
 \leq2\norm{(I-\mathsf E_M)f}_2.
 \label{a-limit-zero-hamiltonian-extension}
\end{equation}
Likewise, for nonreal $z$,
\begin{align}
 (z-\widetilde H_M)^{-1}
 &=\mathsf E_M(z-H_\epsilon)^{-1}
   +z^{-1}(I-\mathsf E_M),\nonumber\\
 \norm{((z-\widetilde H_M)^{-1}-(z-H_\epsilon)^{-1})f}_2
 &\leq\left(|z|^{-1}+|\operatorname{Im}z|^{-1}\right)
 \norm{(I-\mathsf E_M)f}_2.
 \label{a-limit-strong-resolvent}
\end{align}
Thus $\widetilde H_M$ converges to $H_\epsilon$ in the strong resolvent
sense. These statements specify operators on one common Hilbert space.
They do not compare matrices of changing dimensions without an embedding.

The embedded propagators converge strongly, uniformly in time on each
fixed state. Their operator norm error nevertheless satisfies
\[
 \norm{U(t)-U_M^{\mathrm{emb}}(t)}=1
 \qquad\text{for every }M,t.
\]
The fully unitary extensions also converge strongly with the bound in
\eqref{a-limit-zero-hamiltonian-extension}; norm convergence of these
extensions needs a separate condition on the omitted phases. Neither
assertion is a general norm-limit theorem for arbitrary quantum dynamics.

\subsection{A graph-domain example with exact infinite tails}

\begin{example}[Convergent states and finite energy derivatives]
\label{a-limit-geometric-example}
Choose one unit vector $u_m\in\operatorname{Ran}\mathsf F_m$ at every
level, let $0<q<1$, and set
\[
 f_q=\sum_{m\geq0}\sqrt{(1-q)q^m}\,u_m,
 \qquad \epsilon_m=2^m.
\]
Then $\norm{f_q}_2=1$ and
\begin{equation}
 \norm{(I-\mathsf E_M)f_q}_2^2=q^{M+1}.
 \label{a-limit-geometric-state-tail}
\end{equation}
However, $f_q\in\Dom(H_\epsilon)$ exactly when $4q<1$:
\[
 \norm{H_\epsilon f_q}_2^2
 =(1-q)\sum_{m\geq0}(4q)^m.
\]
For $q<1/4$, its graph tail is exactly
\begin{equation}
 q^{M+1}+\frac{(1-q)(4q)^{M+1}}{1-4q}.
 \label{a-limit-geometric-graph-tail}
\end{equation}
For $q=1/8$, the state norm tail decays faster than the graph norm tail.
For $q=1/2$, the normalized state still has a convergent $L^2$ cylinder
approximation and a well-defined unitary evolution, but is outside the
Hamiltonian domain. All these conclusions concern genuinely infinite
series, rather than a finite last level with its tail set to zero.
\end{example}

\adicfigure{10_tail_graph_norm}{Exact infinite tails for the geometric
level masses of Example \ref{a-limit-geometric-example}. The operator
domain controls a weighted tail in addition to the state norm tail.
The parameters are declared in the plate and its reproducible Python
source.}

\subsection{Compact resolvents and thermal weights}

The omitted embedded resolvent has the exact norm
\begin{equation}
 \norm{(I-\mathsf E_M)(z-H_\epsilon)^{-1}}
 =\sup_{m>M}|z-\epsilon_m|^{-1}.
 \label{a-limit-resolvent-tail-norm}
\end{equation}
The upper bound follows by Parseval. For the lower bound, test on a unit
vector at any omitted level and take the supremum. Since every level is
nonzero and finite dimensional, this proves the following equivalence.

\begin{proposition}[Exact compactness criterion]
\label{a-limit-compact-resolvent}
The operator $H_\epsilon$ has compact resolvent if and only if
$|\epsilon_m|\to\infty$. In that case the embedded finite resolvents
converge in operator norm. This is an assertion about the embedded
resolvents, rather than the zero-Hamiltonian extensions of
\eqref{a-limit-strong-resolvent}.
\end{proposition}
\begin{proof}
If the absolute energies tend to infinity, the right side of
\eqref{a-limit-resolvent-tail-norm} tends to zero. The retained part of the
resolvent has finite rank, so the resolvent is a norm limit of finite-rank
operators and is compact. Conversely, if the absolute energies do not
tend to infinity, there are infinitely many distinct levels $m_j$ with
$|\epsilon_{m_j}|\leq C$ for some $C$. Unit vectors $u_j$ on these levels
are orthogonal. For a fixed nonreal $z$, their resolvent images are also
orthogonal and have norm at least $(|z|+C)^{-1}$.
They have no convergent subsequence. The resolvent is therefore not
compact.
\end{proof}

When $|\epsilon_m|\to\infty$, the omitted part of the resolvent for
$\widetilde H_M$ in \eqref{a-limit-strong-resolvent} still carries the
constant multiplier $z^{-1}$. It does not converge in norm to the omitted
resolvent of $H_\epsilon$. Strong resolvent convergence and norm
convergence of an embedded finite-rank resolvent are distinct claims.

For $\beta>0$, the multiplier $e^{-\beta H_\epsilon}$ is bounded precisely
when the energies are bounded below. If there is no lower bound, the
levelwise exponential is an unbounded operator with domain
$\{f:\sum_m e^{-2\beta\epsilon_m}\norm{\mathsf F_mf}_2^2<\infty\}$.
Thus a finite Euclidean intertwining identity does not by itself supply a
bounded infinite heat semigroup.
When bounded below, the positive operator has trace
\begin{equation}
 \operatorname{Tr}(e^{-\beta H_\epsilon})
 =e^{-\beta\epsilon_0}
 +\sum_{m\geq1}(p-1)p^{m-1}e^{-\beta\epsilon_m},
 \label{a-limit-thermal-trace}
\end{equation}
possibly infinite. This follows by summing its nonnegative diagonal
entries in the complete Haar basis. It is trace class exactly when this
series converges. For $\epsilon_m=c m$ with $c>0$, the condition is
$\beta c>\log p$, by the geometric-series test. For
$\epsilon_m=p^{\alpha m}$ with $\alpha>0$, it converges for every
$\beta>0$, since the exponential energy term eventually dominates the
linear-in-$m$ logarithm of the multiplicity. A thermodynamic partition
function therefore needs an energy choice and a summability condition in
addition to Hilbert-space reconstruction.


\section{Compatible interactions across finitely many primes}
\label{a-limit-multiprime}

\subsection{Tensor embeddings and rectangular cutoffs}

Let $\mathcal P$ be a finite prime set and let $\KK_{\mathrm{tor}}$ be a
finite-dimensional complex space for the retained torsion and dyadic sign
registers. Set
\begin{equation}
 \HH_{\mathcal P}^{\infty}
 =\KK_{\mathrm{tor}}\otimes
 \bigotimes_{p\in\mathcal P}L^2(\Z_p,\mu_p).
 \label{a-limit-multiprime-space}
\end{equation}
The corresponding finite register embeds by
$I_{\KK_{\mathrm{tor}}}\otimes\bigotimes_p\mathcal U_{p,h_p}$.
For a multi-depth $\mathbf h=(h_p)$, define
\[
 \mathsf E_{\mathbf h}
 =I_{\KK_{\mathrm{tor}}}\otimes\bigotimes_p\mathsf E_{p,h_p}.
\]
Factorwise projection identities give
\begin{equation}
 \mathsf E_{\mathbf h}\mathsf E_{\mathbf k}
 =\mathsf E_{\min(\mathbf h,\mathbf k)},
 \qquad \mathsf E_{\mathbf h}f\to f
 \quad\text{as every }h_p\to\infty.
 \label{a-limit-rectangular-projections}
\end{equation}
To justify the second statement, first take an elementary tensor. The
difference of the tensor projections and the identity is a finite
telescoping sum, one factor of which tends strongly to zero in each term,
while all other projection factors have norm at most one. Finite sums of
elementary tensors are dense, and the projections are contractions. This
extends convergence to every $f$.

Refinement in distinct coordinates is path independent because
\[
 (V_p\otimes I)(I\otimes V_q)
 =(I\otimes V_q)(V_p\otimes I)=V_p\otimes V_q.
\]
The identity on $\KK_{\mathrm{tor}}$ leaves the separately retained sign
channel untouched by the averaging. This is a finite set of primes.
An infinite tensor product, a restricted adelic product, and a model
including the Archimedean place each require further choices.

Define the joint detail projections
\begin{equation}
 \mathsf F_{\mathbf r}
 =I_{\KK_{\mathrm{tor}}}\otimes
 \bigotimes_{p\in\mathcal P}\mathsf F_{p,r_p},
 \qquad \mathbf r\in\N^{\mathcal P}.
 \label{a-limit-joint-details}
\end{equation}
They are pairwise orthogonal and have finite-dimensional ranges.
Rectangular sums equal $\mathsf E_{\mathbf h}$, so their strong sum is
the identity. This permits interactions between primes while preserving
the exact depth filtration.

\subsection{A complete block-domain theorem}

The scalar weights are one particularly simple choice. The following
editorial extension states the exact additional input for interacting
blocks. It is useful because the torsion operators inside a joint packet
need not commute with one another.

\begin{theorem}[Self-adjoint packet blocks and compatible finite dynamics]
\label{a-limit-block-hamiltonian}
For every $\mathbf r$, let $A_{\mathbf r}$ be a self-adjoint operator on
the finite-dimensional space
$\KK_{\mathbf r}=\operatorname{Ran}\mathsf F_{\mathbf r}$.
On $\HH_{\mathcal P}^{\infty}=\bigoplus_{\mathbf r}\KK_{\mathbf r}$ define
\begin{align}
 \Dom(H_A)&=\left\{f:\sum_{\mathbf r}
 \norm{A_{\mathbf r}\mathsf F_{\mathbf r}f}^2<\infty\right\},
 \nonumber\\
 H_A f&=\sum_{\mathbf r}A_{\mathbf r}\mathsf F_{\mathbf r}f.
 \label{a-limit-block-domain}
\end{align}
Then $H_A$ is self-adjoint, and finite rectangular packet sums form a
core. Every rectangular range reduces $H_A$, and its finite restriction
$H_A^{(\mathbf h)}$ satisfies exact embedding intertwining with every
larger rectangular restriction. For bounded Borel $F$,
\begin{equation}
 F(H_A)f=\sum_{\mathbf r}F(A_{\mathbf r})\mathsf F_{\mathbf r}f,
 \qquad
 F(H_A)\mathsf E_{\mathbf h}\to F(H_A)
 \quad\text{strongly}.
 \label{a-limit-block-calculus}
\end{equation}
In particular,
\begin{equation}
 \norm{e^{-itH_A}(I-\mathsf E_{\mathbf h})f}
 =\norm{(I-\mathsf E_{\mathbf h})f}.
 \label{a-limit-block-unitary-error}
\end{equation}
\end{theorem}
\begin{proof}
Finite rectangular sums belong to the domain and are dense. Symmetry
holds by symmetry on every block. If $g$ belongs to the adjoint domain
and $k=H_A^*g$, testing against an arbitrary vector in
$\KK_{\mathbf r}$ gives
$\mathsf F_{\mathbf r}k=A_{\mathbf r}\mathsf F_{\mathbf r}g$.
Square summability of $k$ is precisely the domain condition in
\eqref{a-limit-block-domain}. Thus $H_A^*=H_A$. Truncating the two
square-summable sequences $f_{\mathbf r}$ and
$A_{\mathbf r}f_{\mathbf r}$ proves the graph-core assertion.

The rectangular projection selects a collection of whole blocks, so it
commutes with $H_A$ on the domain and reduces the operator.
The restriction to each retained block is literally the same matrix at
every subsequent depth. This proves the finite intertwining, and
finite-dimensional spectral calculus on each block gives
\eqref{a-limit-block-calculus}. Boundedness of $F$ bounds its omitted
image by $\norm F_\infty\norm{(I-\mathsf E_{\mathbf h})f}$, proving
strong convergence. Finally, each block exponential is unitary, so
orthogonality of the omitted blocks gives
\eqref{a-limit-block-unitary-error}.
\end{proof}

For example, on a joint depth packet one may take
\[
 A_{\mathbf r}
 =\left(\sum_a\epsilon_{a,\mathbf r}T_a\right)
 \otimes I_{\mathrm{phase},\mathbf r},
 \qquad \epsilon_{a,\mathbf r}\in\R,\quad T_a=T_a^*.
\]
The matrices $T_a$ act on the sign and torsion factors. Their mutual
commutation is unnecessary because their sum is self-adjoint on the
finite block. The coefficients assigned to an existing packet are fixed
under refinement. An extension acting wholly inside a new detail
complement has no matrix element between that complement and the old
range. A cross term between them would be a different model.

There is also an exact construction for any \emph{specified}
self-adjoint compatible finite family along a cofinal nested sequence.
Suppose finite-dimensional spaces $\mathcal V_0\subset\mathcal V_1
\subset\cdots$ have dense union, and self-adjoint $K_M$ satisfy
$K_{M+1}f=K_Mf$ for $f\in\mathcal V_M$.
Self-adjointness implies that $\mathcal V_M^\perp\cap\mathcal V_{M+1}$
is invariant: for $g$ there and $f\in\mathcal V_M$,
$\inner{f}{K_{M+1}g}=\inner{K_{M+1}f}{g}=0$.
The initial space $\mathcal V_0$ and the successive difference spaces
therefore have stable self-adjoint finite blocks.
Their orthogonal direct sum with the maximal square-summability domain
is self-adjoint by the same adjoint proof above. This is an editorial
compatible-family construction, not an assertion that arbitrary finite
Hamiltonians are compatible. The identities $K_{M+1}|_{\mathcal V_M}=K_M$
are its substantive hypothesis.

\subsection{An exact cross-prime entangling generator}

Let $Z_{\mathrm{sign}}$ have eigenvalues $+1,-1$ on a dyadic sign qubit,
and let $B=B^*$ act on the remaining factors. Assume there are orthonormal
states $|b_0\rangle,|b_1\rangle$ with
$B|b_0\rangle=0$ and $B|b_1\rangle=|b_1\rangle$.
One concrete choice is an odd-prime torsion involution times a fixed
odd-prime detail projector, with the torsion state in its $+1$ eigenspace.
For real $g\neq0$, use
\[
 H_{\mathrm{int}}=gZ_{\mathrm{sign}}\otimes B,
 \qquad
 |\Psi(0)\rangle=|+_x\rangle\otimes
 \frac{|b_0\rangle+|b_1\rangle}{\sqrt2}.
\]
Evolution under this explicit interaction alone gives
\begin{equation}
 |\Psi(t)\rangle=\frac1{\sqrt2}
 \left(|+_x\rangle\otimes|b_0\rangle
 +e^{-igtZ_{\mathrm{sign}}}|+_x\rangle\otimes|b_1\rangle\right).
 \label{a-limit-entangled-state}
\end{equation}
Tracing out the orthogonal $b$ states yields the equally weighted mixture
of the two displayed sign states. Their overlap is $\cos(gt)$.
A mixture of two unit vectors with equal weights and overlap magnitude
$c$ has eigenvalues $(1+c)/2$ and $(1-c)/2$: its trace is one and its
determinant on their two-dimensional span is $(1-c^2)/4$.
Thus the sign entropy in bits is
\begin{equation}
 S_{\mathrm{sign}}(t)
 =h_2\left(\frac{1+|\cos(gt)|}{2}\right),
 \quad h_2(u)=-u\log_2u-(1-u)\log_2(1-u),
 \label{a-limit-entanglement-entropy}
\end{equation}
with $0\log_2 0=0$. At $t_* =\pi/(2|g|)$ the two sign states are
orthogonal and the entropy is one bit. A fixed level projector is
compatible under refinement, so this interaction can entangle the prime
factors while preserving the cylinder embedding.
The tensor factorization alone creates no entanglement. It is this
chosen generator and initial superposition that produce the result.


\section{Controlled leakage when compatibility is broken}
\label{a-limit-leakage}

The preceding identities are exact because the coarse range reduces the
Hamiltonian. A weighted coupling to its complement has a different,
quantifiable behavior. This section concerns finite-dimensional
self-adjoint operators; no unbounded perturbation theorem is implicit.

Let $V$ be an isometry from a coarse space to a finite fine space, set
$P=VV^*$ and $Q=I-P$, and let $H_0=H_0^*$ satisfy
$H_0V=VH_c$. Let $B=B^*$ be purely off diagonal:
\[
 PBP=0,\qquad QBQ=0.
\]
For a real $\lambda$, define $H_\lambda=H_0+\lambda B$. In the source
model $\lambda=p^{-w}$ for an assigned weight $w\geq1$.

\begin{proposition}[Defect, leakage, and a second-order self-energy]
\label{a-limit-leakage-bounds}
Under these hypotheses,
\begin{align}
 \norm{H_\lambda V-VH_c}&=|\lambda|\norm{BV},
 \label{a-limit-defect-norm}\\
 \norm{Qe^{-itH_\lambda}P}
 &\leq\min\{1,|t\lambda|\norm{QBP}\}.
 \label{a-limit-duhamel-bound}
\end{align}
For $z\notin\sigma(QH_\lambda Q|_{Q\HH})$, the coarse self-energy is
\begin{equation}
 \Sigma_P(z)=PH_\lambda Q
 (z-QH_\lambda Q|_{Q\HH})^{-1}QH_\lambda P,
 \label{a-limit-feshbach-definition}
\end{equation}
where the inverse acts only on $Q\HH$, and
\begin{equation}
 \norm{\Sigma_P(z)}
 \leq\frac{|\lambda|^2\norm{PBQ}^2}
 {\operatorname{dist}(z,\sigma(QH_\lambda Q|_{Q\HH}))}.
 \label{a-limit-feshbach-bound}
\end{equation}
\end{proposition}
\begin{proof}
The defect is $\lambda BV$, so the first identity is immediate.
Self-adjointness of $H_0$ makes its invariant coarse range reducing.
Hence $Qe^{-itH_0}P=0$, and Duhamel's identity gives
\[
 Qe^{-itH_\lambda}P
 =-i\lambda\int_0^t
 Qe^{-i(t-s)H_\lambda}B e^{-isH_0}P\,ds.
\]
The off-diagonal assumption and reduction by $H_0$ give
$B e^{-isH_0}P=QBP e^{-isH_0}P$. The two propagators and the exterior
projections are contractions. Integrating the norm bound over an
interval of length $|t|$ gives the second expression in
\eqref{a-limit-duhamel-bound}. The bound one follows directly from
unitarity and the norms of the projections.

The coarse/detail blocks of $H_0$ vanish, so
$PH_\lambda Q=\lambda PBQ$ and $QH_\lambda P=\lambda QBP$.
The detail compression is self-adjoint. Diagonalizing it bounds the
resolvent norm by the reciprocal spectral distance. Since
$QBP=(PBQ)^*$, submultiplicativity gives
\eqref{a-limit-feshbach-bound}.
\end{proof}

The small-time leading term is $-it\lambda QBP$. More precisely,
\begin{equation}
 \norm{Qe^{-itH_\lambda}P+it\lambda QBP}
 \leq\frac{|\lambda|t^2}{2}
 \bigl(\norm{H_\lambda}+\norm{H_0}\bigr)\norm{QBP}.
 \label{a-limit-shorttime-remainder}
\end{equation}
To verify this, subtract the constant integrand $QBP$ in Duhamel's
formula and use $\norm{e^{-iuK}-I}\leq|u|\norm K$ for a finite
self-adjoint $K$. The two deviations are bounded by
$(|t-s|\norm{H_\lambda}+|s|\norm{H_0})\norm{QBP}$ on the oriented
integration interval. Their integrals are both $t^2/2$.
For $|\lambda|\leq1$, this is a uniform $O(t^2|\lambda|)$ remainder.
A generic $O(t^2\lambda^2)$ remainder would omit the first-order
interaction with the background Hamiltonian.

For the stipulated assignment $\lambda=p^{-w}$, the operator defect and
this leakage upper bound scale linearly with $p^{-w}$. The self-energy
bound is quadratic in $p^{-w}$ away from the detail spectrum. Under the
purely off-diagonal hypothesis, $QH_\lambda Q=QH_0Q$, so the denominator
does not depend on $\lambda$. Approaching a detail resonance removes
that uniform estimate. These are exact finite quantum statements for a
chosen coupling law. An integral valuation packet does not, by itself,
force every physical interaction strength to equal $p^{-w}$.


\section{A local walk on physical position and refinement depth}
\label{a-limit-product-walk}

\subsection{The finite graph and its unitary tick}

Let $N\geq3$ and $h\geq0$. Let $C_N$ be the physical cycle and let
$T_{p,h}$ be the rooted finite $p$-ary tree of prefixes of length at most
$h$. The product graph is
\[
 \mathcal G=C_N\square T_{p,h}.
\]
A vertex is $v=(x,\mathbf d)$, with $x\in\Z/N\Z$ and $\mathbf d$ a
prefix. An edge changes either $x$ by one on the cycle or the prefix by
one parent-child step. The graph has
$N\sum_{r=0}^h p^r$ vertices. It includes all tree depths, unlike a
cylinder register which contains only the residues at its last depth.

Let $\vec E$ denote the directed arcs and let $\tau_2=2$,
$\tau_p=p-1$ for odd $p$. The finite arc Hilbert space is
\[
 \HH_{\mathrm{arc}}=\ell^2(\vec E)\otimes\C^{\tau_p}.
\]
For each vertex of degree $d_v$, define the normalized outgoing vector
$s_v=d_v^{-1/2}(1,\ldots,1)$ and the Grover coin
$G_v=2|s_v\rangle\langle s_v|-I$.
Its square is $I$, because $|s_v\rangle\langle s_v|$ is a projection.
Take any local unitary $R_v$ on the torsion or sign factor and put
\[
 C=\bigoplus_v(G_v\otimes R_v).
\]
The reversal shift $S$ sends the arc $v\to u$ to $u\to v$, retaining
the internal register. It is a permutation with $S^*=S$ and $S^2=I$.
The one-tick walk
\begin{equation}
 U_{p,h,N}=SC
 \label{a-limit-walk-unitary}
\end{equation}
is unitary because $U_{p,h,N}^*U_{p,h,N}=C^*S^*SC=C^*C=I$.
Here locality concerns the graph edges, with no continuum limit assumed.

\subsection{Local conservation and the product cone}

Write $\psi^n(v,u)$ for the internal amplitude on an arc leaving $v$
before the coin, and $\phi^n=C\psi^n$ for its post-coin value. Define
\[
 \rho_v^n=\sum_{u\sim v}\norm{\psi^n(v,u)}^2,
 \qquad F_{v\to u}^n=\norm{\phi^n(v,u)}^2.
\]
The local unitary coin gives
$\sum_uF_{v\to u}^n=\rho_v^n$.
The reversal shift gives
$\rho_v^{n+1}=\sum_u F_{u\to v}^n$. Subtraction proves
\begin{equation}
 \rho_v^{n+1}-\rho_v^n
 =\sum_{u\sim v}(F_{u\to v}^n-F_{v\to u}^n).
 \label{a-limit-walk-continuity}
\end{equation}
For every set of vertices $A$, summing this equation cancels every
internal directed flux once with its opposite occurrence. Hence
\begin{equation}
 \sum_{v\in A}(\rho_v^{n+1}-\rho_v^n)
 =\sum_{\substack{u\notin A,\ v\in A\\u\sim v}}F_{u\to v}^n
 -\sum_{\substack{v\in A,\ u\notin A\\u\sim v}}F_{v\to u}^n.
 \label{a-limit-walk-coarse-balance}
\end{equation}
Taking $A$ to be all vertices gives conservation of total probability.
The coarse identity uses a subset of vertices and does not presume
amplitude-level Haar refinement.

\begin{proposition}[Exact finite propagation]
\label{a-limit-walk-cone}
If the initial outgoing amplitudes are supported at $v_0$, then
\begin{equation}
 d_{C_N}(x,x_0)+d_T(\mathbf d,\mathbf d_0)>n
 \quad\Longrightarrow\quad\rho_{(x,\mathbf d)}^n=0.
 \label{a-limit-product-cone}
\end{equation}
For two prefixes, writing $\operatorname{lcp}$ for the length of their
longest common prefix,
\[
 d_T(\mathbf d,\mathbf d_0)
 =|\mathbf d|+|\mathbf d_0|-2\operatorname{lcp}(\mathbf d,\mathbf d_0).
\]
In particular, each tick allows at most one physical edge or one tree
edge, and $\bigl||\mathbf d|-|\mathbf d_0|\bigr|\leq n$.
\end{proposition}
\begin{proof}
Every path in a Cartesian product graph uses a number of physical edges
at least the physical distance and a number of tree edges at least the
tree distance. Their sum is a lower bound on path length. Following a
shortest physical path and then a shortest tree path attains that sum,
so it is the product graph distance. In a rooted tree the shortest path
rises to the longest common prefix and then descends, giving the
displayed formula.

The claim at $n=0$ is the support hypothesis. A coin mixes outgoing
arcs only at their common source, so it does not enlarge the vertex
support. The shift moves the source across one edge. A state supported
in the radius-$n$ ball is therefore supported in the radius-$(n+1)$
ball after a tick. Induction proves the result. Since each graph edge
changes only one factor, the individual distance and depth bounds follow.
\end{proof}

The depth difference alone is weaker than the full tree distance.
Different prefixes at the same depth can have positive distance, so a
plot aggregated over depth does not show the complete graph geometry.
The cone is an exact support property of this finite coined model.
A Lorentzian continuum dispersion law or a relativistic limiting field
equation requires a separate scaling analysis.

\adicfigure{07_product_walk}{A finite cycle--tree walk, aggregated over
vertices at the same tree depth. The graph-distance support check uses
the full product graph, including prefix labels. Probability plotted
against depth alone displays only an aggregation of that geometry.}

\FloatBarrier
\Needspace{17\baselineskip}
\subsection{Why growing this tree is not automatically Haar refinement}

\begin{example}[A boundary coin breaks the canonical arc embedding]
\label{a-limit-walk-refinement-counterexample}
Take $h\geq1$ and $N\geq3$, and embed the old arcs of
$C_N\square T_{p,h}$ in $C_N\square T_{p,h+1}$ by identical labels,
with zero amplitude on new arcs. Denote this isometry by $J$.
At an old leaf, the old degree is $d=3$: two cycle neighbors and one
parent. The larger graph adds $p$ children, changing the degree to $d+p$.
For an old outgoing arc basis vector $e_j$, the new Grover coin gives
amplitude $2/(d+p)$ on each of the $p$ new outgoing child arcs. An
internal unitary preserves its internal norm. The subsequent reversal
shift sends these amplitudes to new arcs at the children, outside the old
arc range. Therefore, for a unit internal state $\chi$,
\begin{equation}
 \norm{(I-JJ^*)U_{p,h+1,N}J(e_j\otimes\chi)}
 =\frac{2\sqrt p}{3+p}>0.
 \label{a-limit-walk-refinement-leakage}
\end{equation}
The state $JU_{p,h,N}(e_j\otimes\chi)$ has zero amplitude there, so
$U_{p,h+1,N}J\neq JU_{p,h,N}$.
\end{example}

Each walk remains unitary, locally conservative, and confined to its
product graph cone. A different coin family, embedding, or infinite-tree
construction could be investigated. The canonical old-arc inclusion
fails the compatibility hypothesis of
Theorem \ref{a-limit-block-hamiltonian}. The cylinder limit and the tree
walk are exact models with distinct state spaces.


\section{The boundary of the infinite-depth result}
\label{a-limit-scope}

The exact result comprises complex $L^2$ cylinder reconstruction and
dynamics for declared real level energies or self-adjoint packet blocks.
State convergence uses an $L^2$ tail; graph norm convergence uses an
operator-domain tail. Compact resolvents require energies escaping every
bounded set, and finite thermal traces require multiplicity summability.
Each claim has a distinct hypothesis.

The source programs check finite restrictions of these formulas. Code 45
uses three towers,
\[
 (p,M)=(2,7),\ (3,5),\ (5,3),
\]
and a fixed random seed. Their terminal zero tails do not certify an
infinite domain or compactness; the proofs supply those conclusions.
Code 40 tests specified dyadic/ternary interactions and weighted
off-diagonal couplings. Code 42 tests finite product walks, whose
unitarity does not imply compatibility with the Haar tower.

C15 records finite refinement identities and C16 the conditional
infinite-depth construction. Their archived chapter numbers predate the
current arrangement. Exact members are chapter 31 for multi-prime
dynamics, chapter 33 for product walks, and chapter 36 for the infinite
Haar limit. Prime-polyphase semigroups appear in the spectral part.
No fresh Lean compilation, unrestricted quantum-field completion,
infinite-prime adelic construction, or experimental identification is
claimed here.
